\section{三乗根包絡と中心三次切断}\label{sec:seven}
$q^2+q+1=0$ を固定する。$\mathbf a=(a,b,c)$ に対し
\begin{equation}\label{eq:B}
B_{\mathbf a}=\frac{k\langle t,x,y,z\rangle}
{(yz-qzy-tx-at^2,\ zx-qxz-ty-bt^2,\ xy-qyx-tz-ct^2,\ [t,x],[t,y],[t,z])}
\end{equation}
と置く。$d=(q-1)^{-1}=-(q+2)/3$ とし、
\begin{align}
X_h&=x^3+q^2t^2x,\quad Y_h=y^3+q^2t^2y,\quad Z_h=z^3+q^2t^2z,\\
\Omega_{\mathbf a,h}&=xyz+\frac{q-1}{3}t(x^2+z^2)+\frac{q+2}{3}ty^2
                 +dt^2(ax+qby+cz).\label{eq:Omega}
\end{align}

\begin{remark}[普遍Painlev\'e表示との係数比較]
Chekhov--Mazzocco--Rubtsov \cite[Definition 1.1, (1.5)]{CMR} の量子化変数を
本稿と同じ $q$ とし、$\rho^2=q$、$\lambda=\rho(q^{-1}-q)$ と置く。
同文献の式に
\[
 (X_1,X_2,X_3)=(x,y,z),\qquad
 \epsilon_1=\epsilon_2=\epsilon_3=\lambda^{-1},\qquad
 (\Omega_1^0,\Omega_2^0,\Omega_3^0)=(a,b,c)/(q-1)
\]
を代入して各式に $\rho$ を掛けると、本稿の $t=1$ の三関係を得る。
例えば第一式は $xy-qyx-z-c=0$ になる。
三乗根では $\lambda(q-1)\ne0$ なので、この表示の特殊化は正則である。
ただし同文献のDefinition 1.1が仮定する非根条件は満たされない。
以下の三乗根でのPBW性と中心増大は、表示の一致に加えて直接証明する。
\end{remark}

\begin{proposition}\label{prop:pbw}
全係数で $B_{\mathbf a}$ はPBW基底 $t^\ell x^iy^jz^k$ を持ち、
各fiberはnoetherian quadratic Koszul CY AS4型14641の整域である。
低次数中心は
\[
Z(B_{\mathbf a})_1=kt,\qquad
Z(B_{\mathbf a})_3=k\Omega_{\mathbf a,h}\oplus kX_h\oplus kY_h\oplus kZ_h\oplus kt^3
\]
であり、非被約底を含む底変更と可換する。
\end{proposition}
\begin{proof}
$t=1$ の書換えは
$yx=q^2xy-q^2z-q^2c$、$zx=qxz+y+b$、$zy=q^2yz-q^2x-q^2a$。
唯一の重なり $zyx$ の二還元が全係数で一致する。
Rees斉次化とDiamond lemmaからPBW基底を得る。
活動次数の先頭環は三変数量子多項式環 $S$ の中央多項式拡大 $S[t]$ であり、
整域・noetherian性が従う。正則中央拡大のAS正則性 \cite{LBS} と二次Gr\"obner基底からAS4とKoszul性を得る。
Nakayama自己同型は $t$ を固定し $x_i\mapsto x_i+v_it$ と書ける。
第一関係の差は
$(1-q)t(v_yz+v_zy)+((1-q)v_yv_z-v_x)t^2=0$。
PBW独立性から全 $v_i=0$、よってCYである。

式\eqref{eq:Omega}と三つの $X_h$ の十二交換子は元の関係から零となる。
$S$ の中心単項式の条件は $i\equiv j\equiv k\pmod3$。
活動次数一・二には非定数中心がなく、三次には $x^3,y^3,z^3,xyz$ のみがある。
最高記号を引く帰納法で全低次数中心を得る。
非中心係数を消す際に使う $1-q,1-q^2$ は単元なので、底変更についても同じ議論が有効である。
\end{proof}

再開と検証に必要な全中心式を保持する。$t=1$ における元を $X,Y,Z,\Omega$ とし、
$f=(1-q)/3,e=-(2q+1)/9$ と置くと
\begin{align*}
\ell_a&=((2q+1)a^3+2(q-1)a+3qbc)/9,\quad \ell_b,\ell_c\text{ は巡回置換},\\
K&=((q+1)(a^2+b^2+c^2)+(q+2)abc)/3,\\
L&=((1-q)(a^2b^2+a^2c^2+b^2c^2)+(q+2)(a^2+b^2+c^2))/9.
\end{align*}
全中心は一本の関係
\begin{equation}\label{eq:center}
\begin{split}
0={}&\Omega^3+f\Omega^2-XYZ+e(X^2+Y^2+Z^2)\\
&-(aX+bY+cZ)\Omega+\ell_aX+\ell_bY+\ell_cZ+K\Omega+L
\end{split}
\end{equation}
で与えられる。元のPBW関係でこの恒等式を検証できる。
四生成元が先頭中心を生成することと、$\Omega$ に関するmonic三次式の剰余
$X^iY^jZ^k\Omega^r$ ($0\le r<3$) の相異なる最高記号により、これが全関係である。

$u=[\alpha:\beta:\gamma:\delta:\varepsilon]\in\PP^4$ に対して
\[
G_{\mathbf a,u}=\alpha\Omega_{\mathbf a,h}+\beta X_h+\gamma Y_h+\delta Z_h+\varepsilon t^3,
\qquad R_{\mathbf a,u}=B_{\mathbf a}/(G_{\mathbf a,u})
\]
とする。$\Delta=\alpha^3+27\beta\gamma\delta$、
$T_{\rm ell}=\{\beta\gamma\delta\Delta\ne0\}$ と置く。
境界は、平面三次
\[
E_u:\quad \beta x^3+\gamma y^3+\delta z^3+\alpha q^2xyz=0
\]
の $\phi=\operatorname{diag}(1,q,q^2)$ によるtwisted homogeneous coordinate ringである。
$T_{\rm ell}$ 上では三座標頂点を避ける滑らかな楕円曲線であり、$\phi$ は正確な位数三の平行移動となる。

\begin{assumption}[現行H3原稿から使う入力]\label{ass:recognition}
次の結果を既存の証明への明示的な依存として使う。
(1) $\mathbf a=0$ の切断のH3部分は非空開集合 $\mathcal U\subset\PP^4$ である。
(2) 被約有限型底上、PBW平坦な中央包絡・中央三次商族で、中央境界が滑らかな楕円THCR、
その平行移動が同一の有限位数を持つ場合、H3条件は開条件である。
(3) 中央滑らか楕円境界の場合、H3性は $R[t^{-1}]_0$ の有限大域次元と同値である。
(4) H3のpointed exact同値から反標準環と二次包絡が復元される。
証明の所在は研究ノートの固定参照表に記す。これらの認識結果を以下の代数計算のみから結論しない。
\end{assumption}

入力\ref{ass:recognition}よりH3開集合 $\mathcal V\subset\AA^3\times T_{\rm ell}$ を得る。
零切片は $\mathcal U_{\rm ell}=\mathcal U\cap T_{\rm ell}$ に一致するため、$\mathcal V$ は非空で七次元既約である。
全包絡同型は次数スカラーを除いて二符号変更と巡回置換だけであり、
\begin{equation}\label{eq:N}
\mathcal N=[\mathcal V/\mathfrak A_4]
\end{equation}
は七次元の滑らかな分離DM表示stackである。
実際、中心一次直線を保って $x_i\mapsto\sum_jM_{ij}x_j+v_it$ と書き、
境界のpoint graphを比較すると $M$ は対角・巡回型になる。
関係の $tx_j$ 係数が $v_i=0$ を強制し、残りから
$M=\lambda\operatorname{diag}(\epsilon_i)$、$\epsilon_i^2=1,\prod\epsilon_i=1$ を得る。
従って $a=b=c=0$ は余次元三の正則閉部分stackであり、固定包絡四次元族の非成分性が従う。

\section{固定位数三の全中央PBW正常形}\label{sec:normal}
\begin{theorem}\label{thm:normal}
$S=k\langle x,y,z\rangle/(yz-qzy,\mathrm{cyclic})$ の中央PBW拡大を
\[
r_i=t\sum_jA_{ij}x_j+b_it^2,\qquad [t,x_i]=0
\]
と書く。ここで $A$ は対称とする。
$A_{11}A_{22}A_{33}\ne0$ の開部分は、有限\'{e}taleな係数拡大の後、式\eqref{eq:B}へ正規化される。
従って七次元族は、この固定位数三・中央境界部分の指定開集合の全正常形である。
\end{theorem}
\begin{proof}
ポテンシャルを
\[
\Phi=xyz-qxzy-\tfrac12\sum_{i,j}A_{ij}x_ix_j-\sum_i b_ix_i
\]
と書く。二次以下の巡回語は対称式として扱える。
$x_i=X_i+v_i$ とすると、非対角二次係数を消す唯一の平行移動は
\[
v_1=A_{23}/(1-q),\quad v_2=A_{13}/(1-q),\quad v_3=A_{12}/(1-q)
\]
である。対角係数 $h_i=A_{ii}$ は変わらず、新しい定数は
\[
b'_i=b_i+h_iv_i+(1-q)v_jv_k\qquad(\{i,j,k\}=\{1,2,3\})
\]
となる。これはポテンシャルを直接展開した恒等式である。
次に $X_i=r_iY_i$、$t$ は固定とし、
\[
r_i^2=h_jh_k,\qquad r_1r_2r_3=h_1h_2h_3
\]
を選ぶ。単元の平方根を有限\'{e}taleに添加すれば可能で、選択肢は四つの偶符号である。
$h_ir_i=r_jr_k$ により対角一次係数はすべて1、新定数は $b'_i/(r_jr_k)$ となる。
残余群は前節の全同型計算と一致する。
\end{proof}

対称性は単なる便宜上の十分条件ではない。
巡回三次syzygyを一次補正へ作用させると
$\sum_{ij}A_{ij}[x_i,x_j]=0$ を得る。
$q\ne1$ の量子平面では三つの交換子が独立なので反対称部分は零である。
逆に対称 $A$ はKoszul PBW条件を満たす \cite{BG,EG}。
したがって本節は対称な表示だけを恣意的に選んだ分類ではない。
ただし中央境界を失う変形、非正規境界、対角係数零の閉部分はこの定理に含めない。

\section{三捩れから外へ出る全次数の変形}\label{sec:integrate}
\subsection{低次中心がないことによる持上げ}
\begin{lemma}[中央三次の正則持上げ]\label{lem:central-lift}
$A=k[[s]]$、$K=k((s))$ とする。
$V_A$ を濾過各次数で有限自由なPBW族、$\operatorname{gr}V_A=S_A$ とする。
$S_A$ に中央三次元 $g_A$ があり、$g_0$ の中心持上げ $G_0\in V_0$ が存在すると仮定する。
さらに $F_2Z(V_0)=k$、$\dim_K Z(S_K)_3=1$、
$V_K$ に次数三の非定数中心元が存在するとする。
このとき $g_A$ は $V_A$ の中心元へ持ち上がり、その定数項を固定すれば一意である。
指定した $G_0$ を特殊値に取れる。
\end{lemma}
\begin{proof}
$g_A$ の任意の持上げ $\widetilde g$ を選び、補正を
$F_2V_A/A$ の九つの基底で書く。
三生成元との交換子は有限自由加群間の線型方程式 $M(s)c=b(s)$ となる。
$F_2Z(V_0)=k$ なので $M(0)$ の階数は9であり、ある $9\times9$ minorは単元である。
この九行を使った解 $c$ は $A$ に属する。
一般fiberでは非定数中心元の最高記号が $g_K$ の非零倍であり、
規格化すれば全方程式の解が存在する。九行で解は一意なので、残る行も $K$ 上で零になる。
$A$ は整域であるから全行が $A$ 上で零である。
特殊fiberでも定数を除いた解は一意なので所望の $G_0$ へ戻る。
\end{proof}

\subsection{楕円曲線を固定し、平行移動を動かす}
有限な座標拡大の後、境界曲線を
\[
E_\lambda:\quad x^3+y^3+z^3-3\lambda xyz=0,\qquad \lambda^3\ne1
\]
とする。原点を $(1:-1:0)$、三捩れ点を $p_0=(1:-q:0)$ と取る。
\[
p_s=(1:b(s):s),\qquad 1+b(s)^3+s^3-3\lambda b(s)s=0,\qquad b(0)=-q
\]
を考える。$b$ に関する偏微分は $3q^2\ne0$ なので、$b(s)$ は一意の形式解を持ち、
実際に有限型曲線の\'{e}tale局所座標で定義される。
対応する三変数環は
\[
S_s=k\langle x,y,z\rangle/(yz+b(s)zy+sx^2,\ \mathrm{cyclic})
\]
である。

\begin{proposition}[全次数の包絡・切断変形]\label{prop:formal-integral}
七次元族の滑らかな楕円境界部分の任意の幾何学的点は、
境界の平行移動が三捩れから動く、全次数平坦なCY AS4包絡と正則中央三次商の形式族へ延長できる。
その一般fiberも滑らかな同じ楕円曲線を境界とする。
ここでは一般fiberのH3性はまだ結論しない。
\end{proposition}
\begin{proof}
上記のSklyanin表示は $p_s$ による平行移動の二次包絡である。
multilinear行列の行列式は
$b(s)s\{3\lambda xyz-(x^3+y^3+z^3)\}$。
一般fiberは非退化な楕円AS3、特殊fiberは既知の量子多項式環である。
全次数AS正則性は \cite{IS} と標準の楕円THCR認識を使う。
各次数のHilbert値は両fiberで等しいので、DVR上の有限加群の構造定理により $S_A$ は次数別に自由である。
THCR族への全射の三次核は階数一の束をなし、これを $kg_A$ と書く。
平行移動の場合のNakayama公式から $g_A$ は中心であり、$g_0$ は指定した境界三次式である。

ポテンシャルの二次・一次項を特殊fiberの値に固定してPBW変形を取る。
Etingof--Ginzburg \cite[Theorem 3.3.2(i)]{EG} と正則中央拡大の定理により、
Rees包絡は各fiberでCY AS4、全次数で平坦である。
同定理(ii)から一般fiberに非定数中央三次元が存在する。
一般平行移動は位数一・二・三ではなく、先頭環の低次数中心は定数と唯一の三次直線に限られる。
特殊fiberでは命題\ref{prop:pbw}により $F_2Z=k$。
補題\ref{lem:central-lift}を適用すると、指定した中央三次切断が正則に持ち上がる。
境界三次元の正則性と次数別乗法の単射により、商も全次数で平坦となる。
境界THCRが整域なので $t$-進最高記号から商の整域性も従う。
\end{proof}

\begin{lemma}[超越的中心値のH3性]\label{lem:generic-D}
前命題の包絡と中央三次持上げを中心値 $D$ から独立に定義し、
$D$ をそれらの定義体上超越的に取る。このとき一般fiberの切断 $G+D$ はH3である。
\end{lemma}
\begin{proof}
非斉次化 $V$ は有限大域次元を持つCY3環で、$G$ は非定数中心元である。
$G$ の最高記号が非零の正次数正則元なので $G$ は定義体 $F$ 上超越的である。
$F(D)$ 上の $V/(G+D)$ は、$D=-G$ と同定した中心局所化
$V\otimes_{F[G]}F(G)$ に等しい。
中心局所化は有限大域次元を保つ。
滑らかな中央楕円境界に対する入力\ref{ass:recognition}(3)を適用してH3を得る。
これは任意の数値 $D$ に関する開性や許容条件を主張する議論ではない。
\end{proof}

この構成は七次元族の幾何学的一般点から外へ出るH3一般fiberを与える。
境界データ $(E,\rho)$ の二次元、対称二次・一次ポテンシャルの九係数、
平行移動三次元、共通拡大一方向、中心値一方向を数えると
\[
2+9-3-1+1=8
\]
となる。
$1+b(s)$ が単元なので非対角二次項は平行移動で一意に消せる。
一般の境界の射影自己同型群は有限であり、対角二次係数の一つを共通拡大で規格化すれば、
この計算は単なる冗長パラメータ数でなく、八次元の一般表示族の次元を与える。
全intrinsic族への射を大域的に定義したこととは区別する。

\section{滑らかな中心でのHochschild変形空間}\label{sec:hh}
\begin{assumption}[本節の幾何学的適用範囲]\label{ass:order}
$C$ は滑らかな射影三次曲面、$D\in|-K_C|$ は滑らかな楕円曲線とする。
$\Lambda$ は $C$ 上の有限orderで、$C\setminus D$ ではAzumaya、
$D$ 上では\'{e}tale局所的な完備化が
\[
k[[s,w]]\rtimes\mu_3,\qquad \zeta(s,w)=(s,\zeta w)
\]
とMorita同値であるとする。
七次元族の滑らかな中心を持つ開部分は、既存の中心・境界order定理によりこの範囲に入る。
原点切片には第\ref{sec:exact}節の厳密な滑らかさの証明書を持つ点がある。
\end{assumption}

\begin{lemma}[局所Hochschild層]\label{lem:local-HH}
仮定\ref{ass:order}の下で、$C$ 上の局所Hochschild cohomology層は
\[
\mathcal{HH}^0=\OO_C,\quad
\mathcal{HH}^1=T_C(-\log D),\quad
\mathcal{HH}^2=\bigwedge^2T_C(-\log D),\quad
\mathcal{HH}^{>2}=0
\]
である。
\end{lemma}
\begin{proof}
非分岐部分は局所Morita不変性と通常のHKRである。
分岐点では $A=k[[s,w]]$ の二変数Koszul分解を使い、
$\Ext_{A^e}(A,A\rtimes\mu_3)^{\mu_3}$ を計算する。
恒等元の部分は $\bigwedge T_A$ の不変部分であり、その基底は
\[
1;\qquad \partial_s,\ w\partial_w;\qquad w\partial_s\wedge\partial_w.
\]
$z=w^3$ とすると $w\partial_w=3z\partial_z$ なので、これらは
$\OO_C,T_C(-\log D),\wedge^2T_C(-\log D)$ に対応する。

非恒等元の固定部分は $w=0$ の直線である。
Koszul分解の $w$ 方向を計算すると、寄与は法線表現の行列式を伴う。
この一次元表現への $\mu_3$ 作用は非自明なので、次数一・二の両方で不変部分は零である。
これは \cite[Theorem 1.1]{NS} の固定部分公式の直接計算でもある。
この局所計算はBrauer捩れを消す大域的同値を必要としない。
局所Morita同値の変更は、恒等部分の中心に由来するpolyvector層に恒等に作用する。
それらの同定は非分岐稠密開上で通常の中心の同定に一致し、局所自由層への延長が一意であるから貼り合う。
非恒等部分は局所的に零なので、Brauer捩れから新しい層が生じることもない。
\end{proof}

\begin{theorem}[八次元の無障害な圏変形]\label{thm:HH8}
$\mathcal C=\Perf(C,\Lambda)$ とする。仮定\ref{ass:order}の下で
\[
\HH^0(\mathcal C)=k,\qquad \HH^1(\mathcal C)=0,\qquad
\dim\HH^2(\mathcal C)=8,\qquad \HH^n(\mathcal C)=0\ (n\ge3).
\]
さらに標準的な完全列
\[
0\longrightarrow H^1(C,T_C(-\log D))\longrightarrow\HH^2(\mathcal C)
\longrightarrow H^0(C,\OO_C)\longrightarrow0
\]
があり、左右の次元は7と1である。
従って通常のHochschild複体で制御される圏の形式変形は無障害で、
次数スカラーの高次自己同型を別にすれば、そのminiversal底は八次元形式円板である。
\end{theorem}
\begin{proof}
滑らかな三次曲面について $H^{>0}(\OO_C)=0$、
$h^0(T_C)=0,h^1(T_C)=4,h^2(T_C)=0$。
後二つは $C\subset\PP^3$ の法束列とEuler列、
$h^0(\OO_C(3))=19$、$h^0(T_{\PP^3}|_C)=15$ から得る。
$h^0(T_C)=0$ は、連結自己同型が27本の直線を固定し、
六本を縮約した $\PP^2$ 上で一般位置の六点を固定することからも従う。

対数接束の完全列
\[
0\longrightarrow T_C(-\log D)\longrightarrow T_C
\longrightarrow\OO_D(D)\longrightarrow0
\]
を使う。楕円曲線上の $\OO_D(D)$ の次数は3なので $h^0=3,h^1=0$。
従って対数接束のcohomologyの次元は順に $0,7,0$ である。
また $D=-K_C$ から $\wedge^2T_C(-\log D)=\OO_C$。

局所から大域へのHochschild spectral sequenceを補題\ref{lem:local-HH}に適用する。
非零項は $H^0(\OO_C)$、$H^1(T_C(-\log D))$、
$H^0(\wedge^2T_C(-\log D))$ に限られ、すべての微分は次数の理由で零である。
一般的な大域HKR形式性を別途仮定せず、主張と完全列を得る。
通常の変形理論 \cite{LVdB} では接空間が $\HH^2$、小拡大に対する障害が $\HH^3$ に属する。
後者の消滅から形式的滑らかさ、前者の有限次元性から八次元のminiversal底を得る。
\end{proof}

\begin{corollary}
構造対象 $\OO$ がexceptionalである場合、その持上げの障害と選択の曖昧さ \cite{LowenObject} は
それぞれ $\Ext^2(\OO,\OO)$ と $\Ext^1(\OO,\OO)$ に属し、いずれも零である。
従って構造対象を持つ圏の形式変形でも、接空間と障害の上記の値は変わらない。
\end{corollary}

\begin{remark}
この定理は、全H3変形が自動的に同じ14641二次包絡を持つこと、
逆Serre作用が任意の変形heartを保つこと、全次数の相対反標準環が再構成できることを証明していない。
圏の八次元形式変形と、全H3スタックの八次元開部分を同一視するには、これらの比較が必要である。
\end{remark}

\section{厳密計算と具体的な横断方向}\label{sec:exact}
非斉次化 $t=1$、$a=b=c=0$ とし、
\[
\Phi_0=xyz-qxzy-\tfrac12(x^2+y^2+z^2),\qquad
G_0=2\Omega+X+2Y+3Z+5
\]
を取る。$\Delta=170$、$\beta\gamma\delta=6$。
式\eqref{eq:center}に $\Omega=-(X+2Y+3Z+5)/2$ を代入した三次式を $F$ とすると、
$\mathbb Q(q)$ 上で $(F,F_X,F_Y,F_Z)$ のGr\"obner基底は $(1)$ である。
既存の境界正則性定理と合わせ、この点は滑らかな中心と楕円境界を持つH3点である。

次数一から三の巡回語は20個、次数三以下のPBW語も20個ある。
それぞれ $\dot\Phi,\dot G$ の係数を未知数とし、
変形した関係式の下で $[G_0+\epsilon\dot G,x_i]=0\pmod{\epsilon^2}$ を一次化する。
全105行40列の行列 $M$ を $\mathbb Q(q)$ 上で計算すると
\[
\operatorname{rank}M=18,\quad\dim\ker M=22.
\]
12個のaffine座標変更、ポテンシャルのスカラー、切断のスカラーの14列を明示すると、
すべて核に入り階数14である。
既知の七次元族を加えたspanの階数は21であり、次の方向を加えると22になる：
\[
\dot\Phi=x^3+2y^3+3z^3+2xyz,
\]
\begin{align*}
\dot G={}&(22q+6)x+(5q+9)y+(18q+18)z\\
&-\tfrac43(q+1)(x^2+z^2)-\tfrac43(q+4)y^2+6xy+9q\,xz+18yz.
\end{align*}
従って\emph{このポテンシャル・中央三次incidence}の接空間商は正確に八次元である。
これは全H3変形の被覆を行列の階数だけから推論したものではない。

全次数の明示例として、別の許容候補切断
$2\Omega+X+8Y+27Z+D$ に対し
\begin{equation}\label{eq:explicit-integral}
\Phi_u=(1+2u)xyz-qxzy+u(x^3+8y^3+27z^3)-\tfrac12(x^2+y^2+z^2)
\end{equation}
を取る。Etingof--GinzburgのAppendix 9.1の中心式を
\[
X_{\rm R}=x,\quad Y_{\rm R}=2y,\quad Z_{\rm R}=3z,
\quad q_{\rm R}=q/(1+2u),\quad t_{\rm R}=-18u/(1+2u)
\]
で変換すると、$u=0$ で正則な有理中心三次元を得る。
$x^3$ 係数を1、定数を $D$ と規格化した特殊値は正確に
$2\Omega+X+8Y+27Z+D$ である。
全有理係数は小さいJSONに保持し、原関係からの三つの交換子を
$\mathbb Q(q)(u)$ 上で独立に零へ還元した。
この検算は原典の中心性を前提としない。
$D=5$ と固定した一般 $u$ のH3性をこの計算だけで主張しない。
一般中心値 $D$ については補題\ref{lem:generic-D}が適用される。
